\section{Causal defects and same-memory implementation}\label{sec:transfer}
\subsection{A defect of complete conditional cycles}
A cycle comparison in this section includes the physical preparation, all performed costs and audits, visible reports and flags, actions, duration, and the \emph{next retained label}, conditional on every entering label. We use $\TV(P,Q)=\sup_A|P(A)-Q(A)|$. A comparison of report marginals alone is insufficient. For the next proposition the two endpoints share a program interface; a simulator reduction to this situation is stated afterwards.

\begin{theorem}[Recurrence-sensitive causal risk transfer]\label{thm:defect}
Suppose corresponding complete conditional cycle laws have total variation at most $\varepsilon$, uniformly over models, legal programs and entering labels. Suppose both endpoints obey the same \eqref{eq:tail}, have mean lengths at most $\bar\mu$, and satisfy $h(P),h(P')\le H$. Put
\begin{equation}\label{eq:defectmod}
 b(\varepsilon)=\inf_{T\ge1}\{T\varepsilon+2a(T)\},\quad
 H_t=4\bar\mu(2M+1)H,\quad
 \chi(\varepsilon)=\min\{1,2b(\varepsilon)+4H_t[\varepsilon+b(\varepsilon)]\}.
\end{equation}
Then $|g_\theta(\gam)-g'_\theta(\gam)|\le\chi(\varepsilon)$, and the same bound holds for worst-model values optimized over any one common paired class. A uniform decision-loss error $\omega$ adds at most $\omega$. Under a common bound $\E_i\tau^{1+p}\le K$, one may replace $b(\varepsilon)$ by
\begin{equation}\label{eq:holdermod}
 2K^{1/(1+p)}\varepsilon^{p/(1+p)},\qquad0<p\le1.
\end{equation}
The constants can be reduced using direct bounds for $h(S)$ and $h(S')$ in place of $H_t$.
\end{theorem}
\begin{proof}
Couple the \emph{whole completed marked cycle}, including its length and exit label, maximally. Truncating its reward to the first $T$ calls and its length at $T$ gives bounded functionals in $[0,T]$. Consequently
\[
 \norm{P-P'}_\infty\le2\varepsilon,\qquad
 \norm{r-r'}_\infty,\ \norm{d-d'}_\infty\le b(\varepsilon).
\]
Alternatively on the mismatch event of probability at most $\varepsilon$, the discrepancy of either reward or length is bounded by $\tau+\tau'$. Holder's inequality gives \eqref{eq:holdermod}. Equality of complete marked paths includes equality of their rewards for a common readout; readout perturbations are handled separately.

Since $d_i,d'_i\ge1$ and $r'_i\le d'_i$, direct division yields
\[
 \norm{c-c'}_\infty\le2b(\varepsilon),\qquad
 \norm{S-S'}_\infty\le2\varepsilon+2b(\varepsilon).
\]
For any two stochastic matrices, writing $E=S'-S$ and using the group-inverse identities gives the exact relation
\begin{equation}\label{eq:projectionperturb}
 \Pi_{S'}-\Pi_S=\Pi_{S'}E G_S+G_{S'}E\Pi_S.
\end{equation}
Indeed split the left side as $\Pi_{S'}(I-\Pi_S)-(I-\Pi_{S'})\Pi_S$, and substitute $(I-S)G_S=I-\Pi_S$ and $G_{S'}(I-S')=I-\Pi_{S'}$. As $\Pi_S$ is stochastic, Lemma~\ref{lem:roots} now proves \eqref{eq:defectmod}. This argument does not require equal recurrent-class partitions at the endpoints. Changing readout loss by at most $\omega$ changes the long-run per-call loss by at most $\omega$. Taking supremum and infimum over a paired common class completes the proof.
\end{proof}

\begin{proposition}[Primitive comparison and typed composition]\label{prop:primitive}
For classical complete marked kernels, a per-call operational error at most $\kappa$, and a preparation error at most $\kappa_0$, give a conditional completed-cycle error
\begin{equation}\label{eq:primitive-error}
 \varepsilon\le\min\{1,\kappa_0+\bar\mu\kappa\}.
\end{equation}
The same assertion holds for finite-dimensional quantum instruments when $\kappa$ bounds half the integrated trace-norm error of their unnormalized marked outputs on every input state. The original endpoint supplies the mean length in this one-sided estimate.

A parameter-independent causal simulator must specify its action lift, report transformation, private state, preparation convention, clock, and task readout. If a target observer has $M$ labels and its simulator has $R$ labels, direct execution has at most $MR$ joint labels. To apply Theorem~\ref{thm:defect}, its \emph{joint boundary matrix} must have a stated recurrence bound; an observer bound alone does not imply it. Defects of serial complete-cycle simulations add, and private state counts multiply. One-way simulation gives one-way decision transfer; reverse inequalities require a reverse simulator or an independent lower bound.
\end{proposition}
\begin{proof}
Classically couple the preparations and, while matched, use a maximal coupling of the whole next marked transition, including the next hidden state, report, audit and end flag. Summing the per-call mismatch hazard over the original alive calls gives \eqref{eq:primitive-error}. It is not enough to couple reports and ignore a subsequent hidden-state discrepancy.

Quantum mechanically telescope finite marked words by replacing their first $t-1$ instrument factors by the original factors and their remaining factors by the approximating factors. The unnormalized input to replacement $t$ is therefore the original alive state. Its trace weights the one-call error, and the subsequent trace-preserving completed instrument cannot increase the trace norm on Hermitian differences. Sum these trace masses, then let the truncation increase; the mean original lifetime is the sum of alive masses. Preparation is a separate first replacement. This fixes the orientation of the hybrid argument and uses no normalized posterior denominator.

For the simulator, the joint observer/simulator label determines the next legal operation, so the product-state count is literal. Coupling serial marked outputs gives the additive defect by the union bound; applying the preceding theorem at the resulting recurrence bound gives the risk comparison. Task and action compatibility exclude reading future audits or choosing a true-parameter-dependent row. In general these conclusions concern \emph{target-call} average risk. For source physical-time average risk one must change the holding times in \eqref{eq:classratio}. If every target call costs exactly $c$ source calls with zero loss on the other calls, the scored source-time rate is the target-call rate divided by $c$. For variable conversion times there is a separate ratio in each closed class; a ratio of globally pooled means is not justified.
\end{proof}

\subsection{A support-stable architecture with primitive recurrence bounds}
Fix a finite directed graph $\mathcal G$ on $[M]$. Its reflexive reachability relation permits holds even when $\mathcal G$ has no self-loop. Each nonterminal update is supported on the reflexive reachability relation of $\mathcal G$. The preparation row is the identity. At an ending ordinary report, each outgoing edge of $\mathcal G$ is assigned probability at least a fixed rational $\zeta>0$; remaining mass may use the reachability relation. Assume these constraints are feasible and that from every prepared physical state the first ordinary call ends with probability at least $c_0>0$, uniformly over models and actions. Closed classes of the induced boundary chain are exactly the terminal strongly connected components of $\mathcal G$. The class includes adaptive report rules and random actions; it prescribes no common observer-reset label.

\begin{proposition}[A graph certificate for retained recurrence]\label{prop:graph}
For $M\ge2$, every program in this architecture satisfies
\begin{equation}\label{eq:graphbound}
 P_{ij}\ge c_0\zeta\quad((i,j)\in\mathcal G),\qquad
 h(P)\le H_{\mathcal G}:={4(M-1)\over(c_0\zeta)^{M-1}}.
\end{equation}
The architecture may have several terminal components and periodic components. For $M=1$, take $H_{\mathcal G}=0$.
\end{proposition}
\begin{proof}
On the event of ending at the first ordinary call, the incoming label is still $i$ and the ending row gives each designated edge mass at least $\zeta$. This proves the lower bound. Other within-cycle transitions stay within graph reachability, so cannot create a new terminal component. Select one root in each terminal component and a directed path of at most $M-1$ edges from each vertex to its selected terminal root. The chance of hitting some root in the next $M-1$ boundary steps is at least $(c_0\zeta)^{M-1}$. Applying this bound successively by the Markov property gives an expected hitting time at most $(M-1)/(c_0\zeta)^{M-1}$. The stopped-Poisson construction in Lemma~\ref{lem:roots} gives $h(P)$ at most four times this quantity. Aperiodicity or a common one-step minorizing distribution is not used.
\end{proof}

\subsection{Effective two-sided certification without increasing $M$}
The next result distinguishes an effective input specification from a complexity claim. Assume the raw graph architecture above and \eqref{eq:tail}, with computable bounds. Supply finite executable report partitions, refining the visible types and graph support regions, with complete marked-kernel reconstruction error $\kappa_h\to0$. Reconstruction retains the exact joint cell/next-physical-state/audit/flag output and redraws only the within-cell report from the reference conditional law. Supply finite model nets with a controller-uniform average-risk modulus $\Omega(s)\to0$, a compact decision net of mesh $\rho$ with loss modulus $\omega_D(\rho)$, and certified finite-history integrations and tail truncations. Proposition~\ref{prop:primitive} and Theorem~\ref{thm:defect} provide $\Omega$ when a uniform primitive model modulus is available.

\begin{theorem}[Global lower and implementable upper at the same memory]\label{thm:compile}
Let $V_{\mathcal G}^\Theta(M)$ be the worst-model average value over all Borel programs of the displayed graph architecture. The supplied data yield finite intervals $[L,U]$, with $U-L$ arbitrarily small, and an actual $M$-state table such that
\begin{equation}\label{eq:certificate}
 \max\{0,L-e_{h,Q,\rho}\}\le V_{\mathcal G}^\Theta(M)
 \le\min\{1,U+\Omega(s)\},
\end{equation}
where, using any valid common recurrence bound for the endpoints,
\begin{equation}\label{eq:compileerror}
 e_{h,Q,\rho}=\chi(\bar\mu\kappa_h)
 +\chi\bigl(\bar\mu(|A|+M)/Q\bigr)+\omega_D(\rho).
\end{equation}
For partition sizes $J_a$ and decision-net size $D_\rho$, a deliberately crude enumeration bound is
\begin{equation}\label{eq:count}
 (Q+1)^{M|A|+M^2\sum_aJ_a}\,D_\rho^M,
\end{equation}
with harmless extra finite preparation entries if these are variable. The lower endpoint applies to every Borel program in this architecture, not merely a chosen local optimizer. The upper table operates on the original experiment, not a simulator with additional persistent states.
\end{theorem}
\begin{proof}
Average each Borel report row once on each reference cell. By the definition of reconstruction, running the original row on the reconstructed instrument has exactly the same joint marked law as running this cell-averaged row on the original instrument. This equality includes the retained exit label and holds from every entering label. Convex support and floor constraints are preserved by averaging. Theorem~\ref{thm:defect} and Proposition~\ref{prop:primitive} therefore bound the original-to-cell loss by the first term of \eqref{eq:compileerror}, with the original all-program tail certificate inherited by the reconstructed endpoint through this equality.

For a row having $m$ required edges, reserve the integer mass $Q\zeta$ on each required edge, choosing $Q$ divisible by the denominator of $\zeta$. Distribute the remaining $Q(1-m\zeta)$ integer units by cumulative rounding of the residual probabilities. This preserves all forbidden zeros and required floors, keeps the actual denominator equal to $Q$, and gives row total variation at most $M/Q$. Round action probabilities with error at most $|A|/Q$ and readouts on the decision net. Every rounded table stays in the same graph architecture, with the same uniform tail and recurrence bounds. The cycle coupling gives the second error in \eqref{eq:compileerror}. Thus \emph{every} Borel program is covered by a legal finite table with that common error.

For each finite table and model-net point, truncate physical cycles and integrate the finite words to enclose $P,d,r$. Their tail errors follow from \eqref{eq:tail}. The graph determines the number $k$ of closed classes. For $L_S=I-S$, the identity
\begin{equation}\label{eq:forest}
 \Pi_S={ [z^{k-1}]\operatorname{adj}(zI+L_S)
                   \over [z^k]\det(zI+L_S)}
\end{equation}
follows from $\Pi_S=\lim_{z\downarrow0}z(zI+L_S)^{-1}$ and semisimplicity of zero. The directed forest expansion \cite{ChebotarevAgaev} makes the denominator a sum of nonnegative spanning-forest weights. A forest using the graph edges exists, so it is at least $(c_0\zeta/\bar\mu)^{M-k}$. Hence certified integrations give arbitrarily narrow enclosures for every finite table's $g$ without dividing by an unverified small number. Finite maxima over models and minima over tables give $L$; evaluating the table with least certified upper value gives $U$. Refining intervals makes $U-L$ arbitrarily small. Equation~\eqref{eq:count} bounds the number of integer-coordinate rows before simplex constraints and is intentionally an overcount.

Finite-table inclusion alone does not give the displayed lower bound. The all-Borel coverage just proved gives $V_{\mathcal G}^{F}\ge L-e_{h,Q,\rho}$ on the finite model net $F$; full-family value is no smaller. For the selected finite table, the controller-uniform model modulus gives its full-family risk at most $U+\Omega(s)$. Clipping to $[0,1]$ proves \eqref{eq:certificate}.
\end{proof}

The theorem deliberately uses a support-stable architecture. Rounding an arbitrary member of $\Gamma_{M,H}^\Theta$ can create a new transition between closed classes and need not preserve its recurrence bound. An effective theorem for \emph{all} such spectral sublevels would require another robustness certificate; it is not obtained by silently applying the preceding enumeration.

A table with $P_0=M|A|+M^2\sum_aJ_a$ probability entries uses at most $O(P_0\log(Q+1))$ program bits, in addition to report-partition and decision descriptions. It has only $\lceil\log_2M\rceil$ persistent label bits. A dyadic sampler uses a stated finite number of fresh coin bits per row; if a rational floor is not dyadic, exact rational sampling may use a variable fresh-coin count with finite expectation, or a dyadic smaller certified floor can be used. Cell location, decision evaluation, temporary arithmetic and planning enumeration are separately charged. No bound in \eqref{eq:count} is called efficient, and this implementation does not prove optimal lower bounds for those additional resources.
